The Coma galaxy cluster (Abell 1656) has an angular diameter on the sky of
about 100 arcminutes, and contains more than 1000 individual galaxies, most of
which are dwarf and giant ellipticals orbiting the common center of mass of
the cluster in approximately circular orbits. The table below lists the
measured radial velocities of a few individual cluster member galaxies.

\begin{center}
\begin{tabular}{cc|cc|cc|cc}
No. & $v_r$ (km/s) & No. & $v_r$ (km/s) & No. & $v_r$ (km/s) &
No. & $v_r$ (km/s) \\
\hline
1 & 6001 &  6 & 7116 & 11 & 7156 & 16 & 7111 \\
2 & 7666 &  7 & 7004 & 12 & 7522 & 17 & 8292 \\
3 & 6624 &  8 & 4476 & 13 & 7948 & 18 & 5358 \\
4 & 5952 &  9 & 6954 & 14 & 4951 & 19 & 4957 \\
5 & 5596 & 10 & 8953 & 15 & 7797 & 20 & 7183 \\
\end{tabular}
\end{center}

\begin{description}
\item[a)] Derive the distance of the cluster from the mean radial velocity of
  the galaxies listed in the table. \hfill(8\,p)
\item[b)] Estimate the physical diameter of the cluster (in Mpc).
  \hfill(4\,p)
\item[c)] The virial theorem states that if the galaxy cluster is in dynamic
  equilibrium, then the mean kinetic energy, $\langle K\rangle$, and the mean
  gravitational potential energy, $\langle U\rangle$, are related by
  \[ -2\langle K\rangle = \langle U\rangle \]
  assuming the Coma cluster is spherical. For simplicity, assume that each
  galaxy has approximately the same mass, $m$. Use the virial theorem to prove
  that, in this case, the cluster mass $M$ (also called the virial mass) can
  be expressed as
  \[ M = \frac{5R}{G}\sigma_r^2 \]
  where $\sigma_r^2$ is the velocity dispersion of the cluster. \hfill(10\,p)

  The formula for standard deviation:
  \[ \sigma = \sqrt{\frac{\sum_{i=1}^{n}(x_i - \bar{x})^2}{n-1}} \]
\item[d)] Using the data in the table, estimate the virial mass of the Coma
  cluster in solar masses. \hfill(12\,p)
\item[e)] The total luminosity of the Coma cluster (in solar luminosity,
  $L_\odot$) is $L \approx 5 \times 10^{12}\,L_\odot$. Calculate the
  mass-luminosity ratio of the cluster in solar mass per solar luminosity
  units. \hfill(2\,p)
\item[f)] Which of the following statements is true (more than one answer is
  possible)? Mark the letter of your answer with $\times$ on the answer sheet.
  \hfill(4\,p)
  \begin{description}
  \item[(A)] The mass-luminosity ratio of the Coma cluster is much higher than
    that of a typical spiral galaxy, like the Milky Way.
  \item[(B)] The mass-luminosity ratio of the Coma cluster is similar to that
    of a typical spiral galaxy.
  \item[(C)] The mass-luminosity ratio of the Coma cluster is much less than
    that of a typical spiral galaxy.
  \item[(D)] The Coma cluster contains much more dark matter than a typical
    spiral galaxy.
  \item[(E)] The Coma cluster contains much less dark matter than a typical
    spiral galaxy.
  \end{description}
\end{description}
